\documentclass[11pt]{article}
\usepackage[T1]{fontenc}
\usepackage{lmodern}
\usepackage[margin=0.95in]{geometry}
\usepackage{amsmath,amssymb}
\usepackage{microtype}
\usepackage{enumitem}
\usepackage{xcolor}
\usepackage[colorlinks=true,linkcolor=blue!50!black]{hyperref}

\newcommand{\C}{\mathbb{C}}
\newcommand{\U}{U}
\newcommand{\key}[1]{\par\medskip\noindent{\normalsize\bfseries #1}\par\nobreak\vspace{-1pt}}
\setlist[enumerate]{leftmargin=1.5em,itemsep=3pt,topsep=3pt}

\title{\bfseries Step 5 of 5: one completed execution\\[5pt]
\large\normalfont Companion to \texttt{mub-outline.pdf}, Figure 1; preprint \S7 and Appendix A}
\author{}
\date{}

\begin{document}
\maketitle
\vspace{-2.6em}

\noindent
\textbf{What this step receives.} Four stages of machinery: a finite seed cover, a reduction of each
completion to a point of $\U(2)$, certified cells and balls, and a graph criterion for the last two
bases. Each is an implication from ``a region'' to ``a smaller region'' or to ``no solution''.

\par\medskip\noindent\textbf{What it must produce.} Not another estimate, but a statement about a \emph{run}: that
a documented execution of the whole pipeline visits every region and leaves none undecided. That is
the last thing the proof needs, and it is the one that makes the theorem conditional.

\key{1.\ Records, and the three states}

Every object the earlier stages produce --- a seed box, a retained ball, a refined cell, a vector ball
--- carries a \emph{record} in a finite control flow. A record ends in exactly one of three states:
\begin{enumerate}[label=(\alph*)]
\item \emph{excluded}: a necessary condition failed, with the margin of Appendix A;
\item \emph{resolved}: its completions were transferred to a certified ball which was later excluded;
\item \emph{unresolved}: neither, for instance because a frame check failed or a refinement cap ran out.
\end{enumerate}
The third state is not an error; it is the honest name for ``this region has not been dealt with''. The
difference between a finished proof and an unfinished one is that the list of unresolved records is
empty.

\key{2.\ The implication that is the theorem}

The paper's main statement has the form
\[
\underbrace{\text{a complete execution}}_{\text{a fact about a run}}\ \Longrightarrow\
\underbrace{\text{no four mutually unbiased bases in }\C^6}_{\text{a fact about matrices}} .
\]
Everything a mathematician can check is the implication; the run is the evidence for its hypothesis.
This is worth separating carefully, because two different kinds of claim are being made and they have
different failure modes. A gap in the implication is a mathematical error. A doubt about the run is an
empirical doubt: it would be settled by replaying the pipeline, not by finding a better proof.

\key{3.\ The arithmetic contract}

Each node of the decision tree carries an inequality, and each inequality is evaluated in binary64
floating point. Appendix A is therefore part of the theorem, not an appendix in the editorial sense: it
states the arithmetic and compiler behaviour assumed, provides outward replacements for every operation
used, and proves a comparison margin for every threshold. Without those margins a test that ``passes''
could pass because of rounding, and the exclusion would be worthless. The contract is where the
constants of Steps 1--4 (the $5\pi w^2/4$ of the vanishing sum, the $1.25\pi/N$ of the support bound,
the lowered threshold of the reuse test) are finally converted into statements that survive
floating point.

\key{4.\ The fallback, and why it cannot be skipped}

An unresolved record must be refined: split the region, reprocess the pieces, and try again. The
completed execution is defined as the state in which no record remains unresolved, and the fallback is
what makes that state reachable rather than assumed. Two disciplines go with it. First, refinement has to
control the growth of radii, or the union of retained regions slowly stops containing the solutions it
is supposed to contain. Second, running out of a cap must never be turned into an exclusion --- the
sound move is always to retain the region and admit it.

\key{5.\ What the run is, and what it is not}

The claim is computer-assisted and conditional, and it should be quoted that way. It holds for one
documented execution, under stated binary64 arithmetic and compiler conditions, with the associated
archives identified in the paper's verification directory. The independent checker accompanying the
work does not run the C++ programs, so it does not certify the implementation, and it says nothing
about compiler or hardware behaviour. No specialist peer review of the main claim has been published.
A companion paper proves the weaker bound $N(6)\le5$ by exact arithmetic, and reports of formalisation
concern that weaker bound, not the exclusion of four bases. Nothing above should be read as saying the
upper bound of three has been formalised.

\key{6.\ The shape of the whole argument, in one sentence}

Cut the space into finitely many boxes, map each completion inside a box to a point of a small compact
group, cover that group by cells whose remainders are bounded, cover the last two bases by vector balls,
exclude the clique pattern, and account for every box; then the truth of the accounting, not the truth
of any single estimate, is what closes the problem.

\medskip
\noindent\textit{Note on sources.} The attached excerpt ends in \S3.3. This step is described from the
preprint's \S1.2 summary of the complete execution and from its abstract's statement of Theorem 1.1;
\S7 and Appendix A hold the details.

\end{document}
